% 总图：AI Coach 共用的状态、查证与输出机制。只放图，不放正文段落——
% 正文（「三种角色共用的机制」那段与工具循环分工那段）在 sections/02-delivery.tex 里。
%
% 图号：走 preamble.tex 的 \shotcaption（\refstepcounter{shot}），与 \shot、\shotpair
%       共用同一个计数器，因此全文图号连续、不会重号。正文用 \ref 引用。
%
% 为什么竖排：正文栏只有 453pt 宽，七个框横排需要 500pt 以上，放不下。
%       竖排还有个好处——主链只有「向下」一种走向，顺序不可能读反。
%
% 每框一行：左边步骤号，中间机制名，右边出处文件。步骤号与机制名贴着框的左边界，
%       出处文件贴着右边界——一行两头都有内容，框里不留半边空白，七行等宽、居中、上下对齐。
%       之前每框两行（机制名 + 出处文件）把框高翻了一倍，框又窄又散，七个框看着像浮着的
%       标签，不像一条链。
%
% 行距 0.8cm（框心到框心；框高 16pt，箭头净长 6.7pt）：
%       这是这张图的关键参数。再小箭头顶到框，再大七个框就散成七张标签。
%       整个图连图注约 170pt，不到正文高度（706pt）的 1/3。
%
% 反馈回路不画线：记忆回流到下一轮触发，这件事原来画成一条从最底框拉到最顶框的虚线，
%       那条线与整图等高，读者第一眼看到的是线，不是链。链本身只有「向下」一种走向，
%       顺序已经说明了不回头；回流放在链末一行注里，一句话讲清，不占一行竖线。
%
% 字号：机制名 \footnotesize（10pt），步骤号与出处文件 \scriptsize（8pt），都 ≥7pt，
%       缩到正文栏宽也认得出。出处文件用 \scriptsize 而不是 \tiny（6pt）——\tiny 在这个
%       字号下已经认不清文件名，而文件名是这张图的第二份信息。
\begin{figure}[htbp]\centering
\vspace{4pt}
% 一行的画法：#1 步骤号（同时用作节点名 s#1）、#2 y 坐标、#3 框样式、#4 机制名、#5 出处文件
\newcommand{\prow}[5]{%
  \node[#3] (s#1) at (0,#2) {};
  \node[pnum] at ([xshift=18pt]s#1.west) {#1};
  \node[anchor=west,inner sep=0pt,font=\footnotesize\bfseries]
       at ([xshift=35pt]s#1.west) {#4};
  \node[anchor=east,inner sep=0pt,font=\scriptsize\ttfamily,text=black!65]
       at ([xshift=-15pt]s#1.east) {#5};}
\begin{tikzpicture}[
  font=\scriptsize,
  >={Latex[length=2.2pt,width=2.2pt]},
  % 196pt 装不下整理后的完整路径：实测第 1 行剩 61.3pt 而 src/client/app.js 需 68.0pt、
  % 第 4 行剩 69.7pt 而 src/coach/toolbox.js 需 80.0pt，会压到机制名。
  % 加宽到 232pt 后 7 行全部放得下（本图在 453pt 文字栏内，几何上安全）。
  prow/.style={draw=black!55,rounded corners=2.5pt,minimum width=232pt,minimum height=16pt},
  pterm/.style={prow,fill=black!8},
  pnum/.style={circle,draw=black!45,minimum size=11pt,inner sep=0pt,
               font=\scriptsize,text=black!70},
  flow/.style={->,draw=black!70,line width=0.5pt}]

  % 头尾两框加浅底纹：第一框是入口（谁触发），第七框是出口（记忆落地）
  \prow{1}{0}    {pterm}{玩家提问 / 自动触发}{src/client/app.js}
  \prow{2}{-0.8} {prow} {硬策略门控}          {src/coach/policy.js}
  \prow{3}{-1.6} {prow} {有界工具循环}        {src/coach/runtime.js}
  \prow{4}{-2.4} {prow} {引擎 / RAG 回执}     {src/coach/toolbox.js}
  \prow{5}{-3.2} {prow} {一致性校验}          {src/coach/runtime.js}
  \prow{6}{-4.0} {prow} {有界输出 / 沉默}     {src/coach/runtime.js}
  \prow{7}{-4.8} {pterm}{事件记忆与反思}      {src/coach/memory.js}

  % 主链：七步一路向下，每一步只有一个向下的箭头，六个箭头等长、都落在框的正中
  \foreach \a/\b in {s1/s2,s2/s3,s3/s4,s4/s5,s5/s6,s6/s7} {\draw[flow] (\a) -- (\b);}

  % 回流：不画线，只在链末写一行注
  \node[anchor=north,inner sep=0pt,font=\scriptsize,text=black!65]
       at ([yshift=-9pt]s7.south) {记忆回流到下一次触发（下一轮从第一步重新开始）};
\end{tikzpicture}
\shotcaption{AI Coach 共用的状态、查证与输出机制，以及七个环节各自的出处}\label{fig:pipeline-chain}
\end{figure}
